\section{Appendix} \label{sec:appendix}
In this paper we suggest to calculate the operator from active representations due the feasibility of operations 
\subsection{Passive representation of $\hat{n}_j$ operator}
For this framework we need to specify the initial state, first, for 2 particles in the sites $j$, $j'$ respectively, with $j\neq j'$, the state as function of time is 
\begin{equation} 
|\psi(t) \rangle = e^{-i\mathcal{H} t} |\psi(0)\rangle = e^{-i\mathcal{H} t} \hat{b}_j^\dagger \hat{b}_{j'}^\dagger |0\rangle ,
\end{equation}  
as the temporal propagator is unitary we use the identity to get the following expressions 
\begin{equation}
\begin{split}
    |\psi(t) \rangle &= e^{-i\mathcal{H} t} \hat{b}_j^\dagger \mathcal{I} \hat{b}_{j'}^\dagger \mathcal{I} |0\rangle  \\
         &= e^{-i\mathcal{H} t} \hat{b}_j^\dagger e^{i\mathcal{H} t} e^{-i\mathcal{H} t} \hat{b}_{j'}^\dagger  e^{i\mathcal{H} t} e^{-i\mathcal{H} t} |0\rangle  \\
         &= e^{-i\mathcal{H} t} \hat{b}_j^\dagger e^{i\mathcal{H} t} e^{-i\mathcal{H} t} \hat{b}_{j'}^\dagger  e^{i\mathcal{H} t} |0\rangle,
\end{split}
\end{equation} 
as the bosons obey time invariance $\hat{U}^\dagger \hat{b}^\dagger \hat{U} = \hat{U} \hat{b}^\dagger \hat{U}^\dagger$ 
\begin{equation} 
|\psi(t) \rangle = \hat{b}^\dagger_j \hat{b}_{j'}^\dagger |0\rangle ,
\end{equation}
we derive the Heisenberg equation of motion $\hat{b}^\dagger(t)$ from  
\begin{equation} 
i \frac{d}{dt} \hat{b}_j^\dagger = \left[ \hat{b}_j^\dagger , \hat{H} \right],
\end{equation}
doing the commutator with \eqref{NewHamiBos} obtain: 
\begin{equation} 
i \frac{d}{dt} \hat{b}^\dagger_j(t)= J \left( \hat{b}^\dagger_{j-1} + \hat{b}^\dagger_{j+1} \right) ,
\end{equation} 
consistent with the phenomenology of \eqref{EqEner}, to solve this differential equation we can use the Fourier transform $\hat{b}_j^\dagger = \frac{1}{\sqrt{N}} \sum_k e^{-ikj} \hat{b}_k^\dagger$ and return to the original space with an inverse Fourier transform, so in momentum space the differential equation decouples: 
\begin{equation} 
 i \frac{d}{dt} \hat{b}_k^\dagger (t) = 2J \cos (k) \hat{b}_k^\dagger (t),
\end{equation}
Solving 
\begin{equation} 
 \hat{b}_k^\dagger (t) = e^{-i2Jt\cos (k)} \hat{b}_k^\dagger (0),
\end{equation}
and taking the continuum limit to transform back to time space using the inverse Fourier transform 
\begin{equation} 
\hat{b}_j^\dagger (t) = \frac{1}{2\pi} \int_{-\pi}^{\pi} e^{-ikj} \left( e^{-i2Jt\cos (k)} \hat{b}_k^\dagger (0) \right) dk,
\end{equation}  
substituting $\hat{b}_k^\dagger (0) = \sum_m e^{ikm} \hat{b}_m^\dagger (0)$ 
\begin{equation} 
\hat{b}_j^\dagger (t) = \sum_m \left( \frac{1}{2\pi} \int_{-\pi}^{\pi}  e^{-i2Jt\cos (k)} e^{ik(m-j)} \hat{b}_m^\dagger (0)  dk \right) ,
\end{equation} 
using the Jacobi Anger expansion: 
\begin{equation} 
e^{iz\cos (k)} = \sum_{n=-\infty}^\infty i^n J_n (z) e^{ink}. 
\end{equation} 
Substituting with $z = -2Jt$, we obtain 
\begin{equation} 
\hat{b}_j^\dagger (t) = \sum_{m,n} \left( \frac{1}{2\pi} \int_{-\pi}^{\pi} i^n J_n (-2Jt)  e^{ik(m+n-j)} \hat{b}_m^\dagger (0)  dk \right) , 
\end{equation}
due orthogonality of complex exponentials functions, the integral will be zero unless $n+m-j=0$, so for the values $n= j-m$ and using properties of Bessel function, we get obtain the the expression: 
\begin{equation} 
\hat{b}_j^\dagger(t) = \sum_{m=-\infty}^\infty i^{m-j} J_{m-j}(-2Jt) \hat{b}_m^\dagger(0)
\end{equation}
%\begin{equation}
%\begin{split}
%    \hat{b}_j^\dagger (t) &= \frac{1}{2\pi} \int_{-\pi}^{\pi} e^{ikj} \left( e^{i2Jt\cos (k)} \hat{b}_k^\dagger (0) \right) dk  \\
%         &= \left( \frac{1}{2\pi} \int_{-\pi}^{\pi} e^{ikj}  e^{i2Jt\cos (k)} \hat{b}_k^\dagger (0) dk \right) ,
%\end{split}
%\end{equation} 
\subsection{Active representation of $\hat{n}_j$ operator} 
We obtain for a $|\psi\rangle$ as steady state and $\hat{b}^\dagger$ as dynamic operator, write as: 
\begin{equation} 
\hat{b}_j^\dagger(t) = \sum_{m=-\infty}^\infty i^{m-j} J_{m-j}(-2Jt) \hat{b}_m^\dagger(0), 
\end{equation} 
so the expected value of $\hat{n}_j$ is: 
\begin{equation}
\begin{split}
    \langle \hat{n}_j (t) \rangle &= \langle \psi(0)| \hat{b}_j^\dagger(t) \hat{b}_j (t) | \psi(0) \rangle, \\
         &= \langle \psi(0)| \sum_{m,m'=-\infty}^\infty i^{m-m'} J_{m-j}(-2Jt) J_{m'-j} (-2Jt) \\
         &\quad \hat{b}_m^\dagger   \hat{b}_m | \psi(0) \rangle,
\end{split}
\end{equation}
In general we get a 1-d chain, if we annihilate at $m'$ site and create at $m$ we will obtain an orthogonal state, so for the different values of $m\neq m'$ we will obtain zero, so for an arbitrary initial state we get: 
\begin{equation} 
\langle \hat{n}_j (t) \rangle = \sum_{m=-\infty}^\infty |J_{m-j}(-2Jt)|^2 \langle \hat{n}_m(0) \rangle
\end{equation} 
In general for the active operators are defined with the temporal propagator
\begin{equation} 
\langle \hat{O}(t) \rangle_{+\theta} = \langle \Psi_0 \vert{} e^{i\hat{H}_{B,+\theta} t}  \hat{O} e^{-i\hat{H}_{B,+\theta} t} \vert{} \Psi_0 \rangle
\end{equation}
and it is useful to write the $\mathcal{K}$ invariance defined in \cite{PhysRevLett.121.250404} to obtain some dynamical properties of the observable: 	
\begin{equation}
\begin{split}
\langle \hat{O}(t) \rangle_{+\theta} &= \langle \Psi_0 \vert{} \mathcal{K}^\dagger \left( \mathcal{K} e^{i\hat{H}_{B,+\theta} t} \mathcal{K}^\dagger \right) \\
&\quad \left( \mathcal{K} \hat{O} \mathcal{K}^\dagger \right) \left( \mathcal{K} e^{-i\hat{H}_{B,+\theta} t} \mathcal{K}^\dagger \right) \\
&\quad  \left( \mathcal{K} e^{-i\hat{H}_{B,+\theta} t} \mathcal{K}^\dagger \right) \mathcal{K} \vert{} \Psi_0 \rangle 
\end{split}
\end{equation} 
One should pair $\mathcal{K}$ with $\mathcal{T}$ or $\mathcal{P}$ to restores forward time evolution at angle $-\theta$ or $-U$ respectively, obtaining the following dynamical properties: 
\begin{table}[htpb]
\centering
\caption{Dynamical Relations via $\mathcal{T}\mathcal{K}$ and $\mathcal{P}\mathcal{K}$}
\label{tab:dynamical_relations}
\begin{tabular}{p{0.46\linewidth} p{0.46\linewidth}}
\hline\hline
\textbf{via $\mathcal{T}\mathcal{K}$ ($\theta \leftrightarrow -\theta$)} & \textbf{via $\mathcal{P}\mathcal{K}$ ($U \leftrightarrow -U$)} \\
\hline
\multicolumn{2}{l}{\textit{Site Density $\hat{n}_j$}} \\
$\langle \hat{n}_j(t) \rangle_{+\theta} = \langle \hat{n}_{j'}(t) \rangle_{-\theta}$ & $\langle \hat{n}_j(t) \rangle_{+U} = \langle \hat{n}_{j'}(t) \rangle_{-U}$ \\[1.5ex]

\multicolumn{2}{l}{\textit{Anyonic Bond Energy $\hat{K}_{j,j+1}^{\text{anyon}}$}} \\
$\langle \hat{K}_{j,j+1}^{\text{anyon}}(t) \rangle_{+\theta} $ \newline $= \langle \hat{K}_{j'-1,j'}^{\text{anyon}}(t) \rangle_{-\theta}$ & $\langle \hat{K}_{j,j+1}^{\text{anyon}}(t) \rangle_{+U} $ \newline $= -\langle \hat{K}_{j'-1,j'}^{\text{anyon}}(t) \rangle_{-U}$ \\[1.5ex]

\multicolumn{2}{l}{\textit{Anyonic Current $\hat{I}_{j,j+1}^{\text{anyon}}$}} \\
$\langle \hat{I}_{j,j+1}^{\text{anyon}}(t) \rangle_{+\theta} $ \newline $= -\langle \hat{I}_{j'-1,j'}^{\text{anyon}}(t) \rangle_{-\theta}$ & $\langle \hat{I}_{j,j+1}^{\text{anyon}}(t) \rangle_{+U} $ \newline $= -\langle \hat{I}_{j'-1,j'}^{\text{anyon}}(t) \rangle_{-U}$ \\[1.5ex]

\multicolumn{2}{l}{\textit{Annihilation Operator $\hat{b}_j$}} \\
Couples $\hat{b}_j(t)$ to \newline phase-dressed $\hat{b}_{j'}(t)$ & Couples $\hat{b}_j(t)$ to \newline $(-1)^{j'} e^{i\theta\hat{n}_{j'}} \hat{b}_{j'}(t)$ \\
\hline\hline
\end{tabular}
\end{table}

in the case of anyons using the equation \eqref{NewHamiBos}, .... reducing for $N=2$
\begin{equation}
\begin{split}
\textit{i} \frac{d}{dt} \hat{b}_j(t) &= -J \left( \hat{b}_{j-1} + \hat{b}_{j+1} \right) \\
&\quad - J \Bigl[ \left( e^{-\textit{i} \theta} - 1 \right)\hat{n}_{j-1} \hat{b}_{j-1} \\
&\quad + 2 \left( e^{\textit{i} \theta} - 1 \right) \hat{n}_j \hat{b}_{j+1} \\
&\quad + \left( e^{-\textit{i} \theta} - 1 \right) \hat{b}_{j+1} \hat{b}_j^2 \Bigr]
\end{split}
\end{equation}
%\begin{equation}
%\begin{split}
%    \langle \hat{n}_j (t) \rangle &= \langle \psi(0)| \hat{b}_j^\dagger(t) \hat{b}_j (t) | \psi(0) \rangle, \\
%         &= \langle \psi(0)| \sum_{m,m'=-\infty}^\infty i^{m-m'} J_{m-j}(-2Jt) J_{m'-j} (-2Jt) \hat{b}_m^\dagger   \hat{b}_m | \psi(0) \rangle,
%\end{split}
%\end{equation}
